\chapter{Reading an MRI}

\label{ch:read_mri}

Magnetic Resonance Imaging (MRI) offers unparalleled soft tissue contrast. Unlike CT, which relies on a single tissue property (X-ray attenuation), MRI utilizes magnetic fields and radiofrequency pulses to probe hydrogen proton environments. Interpreting an MRI requires a firm grasp of the different imaging sequences and a systematic reading framework.


\section*{1. MRI Sequences: Deciphering the Contrast}
A clinical MRI scan consists of multiple "sequences," each sensitive to different physical tissue properties. Identifying the sequence is the first step in reading any slice:

\begin{center}
\begin{longtable}{@{}p{0.22\textwidth}p{0.15\textwidth}p{0.18\textwidth}p{0.35\textwidth}@{}}
\toprule
\textbf{Sequence} & \textbf{Water / CSF} & \textbf{Fat} & \textbf{Clinical Utility} \\
\midrule
\textbf{ADC Map} & Dark (black) & Intermediate & Confirms restricted diffusion (matches bright DWI). \\
\textbf{DWI} & Bright (white) & Intermediate & Shows restricted water movement. Highlights hyperacute stroke. \\
\textbf{FLAIR} & Dark (black) & Intermediate-Grey & T2-weighted with free water suppressed. Highlights brain lesions. \\
\textbf{GRE / SWI} & Dark (black) & Intermediate & Susceptibility weighted. Highlights blood products (black). \\
\textbf{T1-Weighted (T1W)} & Dark (black) & Bright (white) & Excellent for anatomical detail. \\
\textbf{T1 + Contrast (Gd)} & Dark (black) & Bright (white) & Highlights blood-brain barrier breakdown (tumors, abscesses). \\
\textbf{T2-Weighted (T2W)} & Bright (white) & Intermediate-Grey & Sensitive to pathology (fluid, edema, inflammation). \\
\bottomrule
\end{longtable}
\end{center}

\subsection*{Key Diagnostic Sequences Explained}
\begin{itemize}
    \item \textbf{FLAIR (Fluid Attenuation Inversion Recovery)}: By suppressing much of the free-CSF signal, FLAIR can make parenchymal lesions more conspicuous. It is useful for assessing demyelinating disease and gliosis; it can also show subarachnoid blood in selected settings, but is not the standalone reference test for acute subarachnoid hemorrhage.
    \item \textbf{DWI (Diffusion-Weighted Imaging) \& ADC (Apparent Diffusion Coefficient)}: In acute ischemia, the failure of sodium-potassium pumps causes cell swelling (cytotoxic edema), trapping water molecules inside cells. This restricted diffusion appears \textbf{bright on DWI} and \textbf{dark on the ADC map}. This combination is highly sensitive for \textbf{ischemic stroke} within minutes of onset.
    \item \textbf{SWI (Susceptibility-Weighted Imaging)}: SWI is sensitive to local magnetic field distortions caused by blood breakdown products (hemosiderin, deoxyhemoglobin) and calcium. These appear black (known as "blooming"). It is essential for detecting microbleeds (amyloid angiopathy), cavernous malformations, and hemorrhagic conversion of ischemic strokes.
\end{itemize}


\section*{2. Reading a Brain MRI: Systematic Approach}
When reading a brain MRI, review the sequences systematically from anatomy to micro-pathology:

\begin{itemize}
    \item \textbf{Alignment and Anatomy (T1W)}: Check sagittal views for cerebellar tonsil position (looking for Chiari malformations) and midline structure alignment. Check axial views for gyral anatomy and midline shift.
    \item \textbf{Brain Parenchyma (T2W \& FLAIR)}: Inspect white matter for hyperintensities (bright white spots). Look for periventricular plaques oriented perpendicular to the lateral ventricles (\textbf{Dawson's fingers} in multiple sclerosis), diffuse microvascular ischemic disease, or focal tumors.
    \item \textbf{Ischemia Survey (DWI/ADC)}: Always scan the DWI and ADC maps together. A matched bright-DWI and dark-ADC lesion represents acute cytotoxic edema (infarction).
    \item \textbf{Contrast Enhancement (T1 + Gadolinium)}: Look for abnormal contrast enhancement:
    \begin{itemize}
        \item \textit{Ring-enhancing lesions}: Differential includes GBM (glioblastoma), metastases, abscesses, or demyelinating lesions.
        \item \textit{Leptomeningeal enhancement}: Suggests meningitis or carcinomatosis.
    \end{itemize}
    \item \textbf{CSF Spaces}: Check lateral, third, and fourth ventricles for enlargement or effacement. Inspect the basal cisterns and sulci.
\end{itemize}


\section*{3. Reading a Joint MRI (MSK MRI)}
MSK MRI requires a systematic evaluation of bone marrow, cartilage, ligaments, and tendons:

\begin{itemize}
    \item \textbf{Bone Marrow (T1W \& STIR/T2-Fat Suppressed)}: Normal marrow has high fat content (bright on T1). Pathological states (fractures, bone bruises, infections) replace fat with fluid, appearing \textbf{dark on T1} and \textbf{bright on STIR/T2-FS} (bone marrow edema).
    \item \textbf{Tendons \& Ligaments}: Healthy tendons and ligaments are composed of dense collagen and contain few protons, appearing as solid \textbf{black bands} on all sequences. A tear presents as a disruption of this black band, filled with bright T2 fluid/edema.
    \item \textbf{Cartilage \& Menisci}: Examine articular cartilage for thinning or focal defects. On knee MRI, inspect the fibrocartilaginous menisci (which appear as black triangles). A tear is represented by a linear high T2 signal that extends to the articular surface of the meniscus.
\end{itemize}


\section*{4. MRI Safety and Contraindications}
Due to the powerful magnetic field (typically 1.5T or 3.0T), safety screening is paramount. 

\begin{itemize}
    \item \textbf{Absolute Contraindications}: MRI safety depends on the exact device, implant, field strength, and conditions of use; complete current MR screening and final-check procedures.
    \item \textbf{Contrast Precautions}: Gadolinium agent choice depends on the indication and current agent-specific guidance. Group II agents have a very low risk of nephrogenic systemic fibrosis (NSF) and should not automatically be withheld solely because eGFR is below 30 mL/min/1.73 m$^2$; assess acute kidney injury and follow local policy.
\end{itemize}


\section*{5. Structured Reporting Example: Multiple Sclerosis}

\begin{testbox}[Indications:]
32-year-old female presenting with optic neuritis and lower extremity paresthesias. Rule out Multiple Sclerosis.
\end{testbox}
\begin{refbox}[Findings:]
\textbf{Brain Parenchyma}: Multiple oval-shaped, T2 and FLAIR hyperintense lesions are noted within the periventricular, juxtacortical, and infratentorial white matter. Several periventricular lesions are oriented perpendicular to the margins of the lateral ventricles, consistent with Dawson's fingers. There is a prominent lesion in the corpus callosum. No significant mass effect or midline shift.\\
\textbf{DWI/ADC}: No areas of acute restricted diffusion to suggest acute ischemic stroke.\\
\textbf{Contrast Enhancement}: Following the administration of intravenous Gadolinium, three of the periventricular lesions and one juxtacortical lesion demonstrate nodular and incomplete ring enhancement, indicating active demyelinating plaques with blood-brain barrier breakdown.\\
\textbf{CSF Spaces \& Ventricles}: Ventricles and sulci are normal for age. No hydrocephalus.\\
\textbf{Bones \& Orbits}: The orbits are normal; no focal abnormalities in the optic nerves. Cranial vault is intact.
\end{refbox}
\begin{pearlbox}[Impression:]
Multiple T2/FLAIR hyperintense white matter lesions satisfying the McDonald criteria for dissemination in space. Active contrast-enhancing lesions confirm dissemination in time. Findings are highly characteristic of active Multiple Sclerosis.
\end{pearlbox}

\clearpage
